% GENERATED FILE. Edit canonical lessons, then run npm run generate:books.
% canonical-source-hash: fnv1a64:c8042a557213a556
% canonical-lessons: GE-C263-nachdem, GE-C263-bevor, GE-C263-sobald, GE-C263-solange, GE-C263-seitdem

\chapter{Nachdem und bevor}
\label{ch:ge-nachdem-und-bevor}
\hlchaptermodality{263}

\begin{chapteropening}
\textbf{By the end of this chapter:} \emph{I can join two past events with after, before, as soon as, as long as and since, with the verb at the end.}

Complete the last lesson of chapter 263: I can join two past events with after, before, as soon as, as long as and since, with the verb at the end.
\end{chapteropening}

\section[nachdem]{nachdem — after (joining two clauses)}

\label{lesson:GE-C263-nachdem}



\begin{quote}
Before the new one: say the German for last of all, then the German for the order.
\end{quote}

\subsection*{You'll want to know: nachdem}
\textbf{nachdem} — \textquotedblleft{}after (joining two clauses)\textquotedblright{}.

\begin{quote}
Nachdem wir gegessen hatten, gingen wir spazieren.
\end{quote}

\begin{quote}
\emph{After we had eaten, we went for a walk.}
\end{quote}

The verb goes to the end, and the earlier event takes the \emph{Plusquamperfekt}: \emph{nachdem er angekommen war}.

\begin{cousinweb}
\emph{nachdem} joins the preposition \emph{nach} \textquotedblleft{}after\textquotedblright{} and \emph{dem} \textquotedblleft{}that\textquotedblright{}, so \textquotedblleft{}after that\textquotedblright{}. Both pieces are inherited Germanic words; EtymDB lists no parent for the compound.
\end{cousinweb}

\begin{grammarlens}[title={What you've built}]
One more word for telling, explaining and describing at B1.
\end{grammarlens}

\subsection*{Your turn}
\begin{itemize}
\raggedright
  \item \emph{Say it:} \emph{nachdem}
  \item \emph{Say it:} \emph{nachdem}, once more
  \item \emph{Say it:} a sentence of your own with \emph{nachdem}
  \item \emph{Recall:} say the German for last of all, then the German for the order, then say \emph{nachdem} again
\end{itemize}

\begin{tcolorbox}[breakable,colback=teal!4,colframe=teal!35!black,title={Before you move on}]
What does nachdem mean? (\textquotedblleft{}After (joining two clauses)\textquotedblright{}.) Say it once more.
\end{tcolorbox}
\section[bevor]{bevor — before (joining two clauses)}

\label{lesson:GE-C263-bevor}



\begin{quote}
Before the new one: say the German for the order, then the German for after.
\end{quote}

\subsection*{You'll want to know: bevor}
\textbf{bevor} — \textquotedblleft{}before (joining two clauses)\textquotedblright{}.

\begin{quote}
Bevor er schlafen ging, las er ein Buch.
\end{quote}

\begin{quote}
\emph{Before he went to bed, he read a book.}
\end{quote}

The verb goes to the end: \emph{bevor der Film beginnt}. The preposition is \emph{vor}: \emph{vor dem Film}.

\begin{cousinweb}
\emph{bevor} is built on \emph{vor} \textquotedblleft{}before\textquotedblright{}, an inherited Germanic word and a cousin of English \emph{fore}, with the old prefix \emph{be-} in front; EtymDB lists no parent for the conjunction.
\end{cousinweb}

\begin{grammarlens}[title={What you've built}]
One more word for telling, explaining and describing at B1.
\end{grammarlens}

\subsection*{Your turn}
\begin{itemize}
\raggedright
  \item \emph{Say it:} \emph{bevor}
  \item \emph{Say it:} \emph{bevor}, once more
  \item \emph{Say it:} a sentence of your own with \emph{bevor}
  \item \emph{Recall:} say the German for the order, then the German for after, then say \emph{bevor} again
\end{itemize}

\begin{tcolorbox}[breakable,colback=teal!4,colframe=teal!35!black,title={Before you move on}]
What does bevor mean? (\textquotedblleft{}Before (joining two clauses)\textquotedblright{}.) Say it once more.
\end{tcolorbox}
\section[sobald]{sobald — as soon as}

\label{lesson:GE-C263-sobald}



\begin{quote}
Before the new one: say the German for after, then the German for before.
\end{quote}

\subsection*{You'll want to know: sobald}
\textbf{sobald} — \textquotedblleft{}as soon as\textquotedblright{}.

\begin{quote}
Sobald der Regen aufhörte, spielten die Kinder im Garten.
\end{quote}

\begin{quote}
\emph{As soon as the rain stopped, the children played in the garden.}
\end{quote}

The verb goes to the end. For the future the present is enough: \emph{Sobald ich Zeit habe, rufe ich an}.

\begin{cousinweb}
\emph{sobald} joins \emph{so} and \emph{bald} \textquotedblleft{}soon\textquotedblright{}. \emph{Bald} is inherited and is a cousin of English \emph{bold} (EtymDB), from an older sense \textquotedblleft{}bold, quick\textquotedblright{}.
\end{cousinweb}

\begin{grammarlens}[title={What you've built}]
One more word for telling, explaining and describing at B1.
\end{grammarlens}

\subsection*{Your turn}
\begin{itemize}
\raggedright
  \item \emph{Say it:} \emph{sobald}
  \item \emph{Say it:} \emph{sobald}, once more
  \item \emph{Say it:} a sentence of your own with \emph{sobald}
  \item \emph{Recall:} say the German for after, then the German for before, then say \emph{sobald} again
\end{itemize}

\begin{tcolorbox}[breakable,colback=teal!4,colframe=teal!35!black,title={Before you move on}]
What does sobald mean? (\textquotedblleft{}As soon as\textquotedblright{}.) Say it once more.
\end{tcolorbox}
\section[solange]{solange — as long as}

\label{lesson:GE-C263-solange}



\begin{quote}
Before the new one: say the German for before, then the German for as soon as.
\end{quote}

\subsection*{You'll want to know: solange}
\textbf{solange} — \textquotedblleft{}as long as\textquotedblright{}.

\begin{quote}
Solange es regnete, blieben wir im Haus.
\end{quote}

\begin{quote}
\emph{As long as it was raining, we stayed in the house.}
\end{quote}

The verb goes to the end. It can also set a condition: \emph{solange du pünktlich bist}.

\begin{cousinweb}
\emph{solange} joins \emph{so} and \emph{lange} \textquotedblleft{}for a long time\textquotedblright{}, built on \emph{lang} \textquotedblleft{}long\textquotedblright{}, an inherited Germanic word (EtymDB).
\end{cousinweb}

\begin{grammarlens}[title={What you've built}]
One more word for telling, explaining and describing at B1.
\end{grammarlens}

\subsection*{Your turn}
\begin{itemize}
\raggedright
  \item \emph{Say it:} \emph{solange}
  \item \emph{Say it:} \emph{solange}, once more
  \item \emph{Say it:} a sentence of your own with \emph{solange}
  \item \emph{Recall:} say the German for before, then the German for as soon as, then say \emph{solange} again
\end{itemize}

\begin{tcolorbox}[breakable,colback=teal!4,colframe=teal!35!black,title={Before you move on}]
What does solange mean? (\textquotedblleft{}As long as\textquotedblright{}.) Say it once more.
\end{tcolorbox}
\section[seitdem]{seitdem — since then; since (joining two clauses)}

\label{lesson:GE-C263-seitdem}



\begin{quote}
Before the new one: say the German for as soon as, then the German for as long as.
\end{quote}

\subsection*{You'll want to know: seitdem}
\textbf{seitdem} — \textquotedblleft{}since then; since (joining two clauses)\textquotedblright{}.

\begin{quote}
Seitdem sie ein Auto hat, fährt sie nicht mehr mit dem Bus.
\end{quote}

\begin{quote}
\emph{Since she has had a car, she no longer goes by bus.}
\end{quote}

Two uses: an adverb, with the verb second (\emph{Seitdem ist er glücklich}), and a conjunction, with the verb at the end. The present covers something still true.

\begin{cousinweb}
\emph{seitdem} joins \emph{seit} \textquotedblleft{}since\textquotedblright{} and \emph{dem} \textquotedblleft{}that\textquotedblright{}. \emph{Seit} comes from Middle High German \emph{sīt} (EtymDB), usually traced to Old High German \emph{sīd} \textquotedblleft{}later, since\textquotedblright{}, a cousin of English \emph{since}.
\end{cousinweb}

\begin{grammarlens}[title={What you've built}]
One more word for telling, explaining and describing at B1.
\end{grammarlens}

\subsection*{Your turn}
\begin{itemize}
\raggedright
  \item \emph{Say it:} \emph{seitdem}
  \item \emph{Say it:} \emph{seitdem}, once more
  \item \emph{Say it:} a sentence of your own with \emph{seitdem}
  \item \emph{Recall:} say the German for as soon as, then the German for as long as, then say \emph{seitdem} again
\end{itemize}

\begin{tcolorbox}[breakable,colback=teal!4,colframe=teal!35!black,title={Before you move on}]
What does seitdem mean? (\textquotedblleft{}Since then; since (joining two clauses)\textquotedblright{}.) Say it once more.
\end{tcolorbox}
